%% n8b_eval.tex -- evaluation, middle part: E4 (operational knowledge and
%% automated distillation) and E5 (content-layer configuration comparison).
%% 2026-11-13: the former seven-configuration tool-layer table (tab:retrieval,
%% fig:retrieval) was removed with the RAG-lane comparison; the verdict-role
%% diagnostic now lives only in E8. Labels owned here: sec:e3, sec:ladder-opt,
%% fig:ladder, fig:miss-overlap.
\subsection{E4: Operational knowledge and automated distillation}
\label{sec:e3}

The experience store holds operational lessons as retrievable
documents---addressed, indexed and searched as an ordinary document, with a
credibility tier deciding whether it may be shown at all
(Section~\ref{sec:org-experience}). Three case measurements exist: an
extraction campaign over one small knowledge base produced thirteen
candidate drafts (archived \path{experience_smoke.json} in the 2026-10-13
batch tarball; \texttt{extract\_candidates}~13, \texttt{drafts\_created}~13);
one end-to-end scenario was sedimented into one experience
(\texttt{exp-d6f4d47ae3d6}) that a later global experience search returned
first, at local similarity 0.711 (archived scenario A6 of the 2026-10-12
full-test campaign); and the draft-approval path was never exercised, because
the draft pool was empty (\texttt{count}~0) and \path{experience_draft_read}
and \path{experience_draft_approve} could not run; the heuristic extraction
path had run only in dry-run mode, returning a dump the quality gate would
reject. One scenario and thirteen candidates cannot support an effectiveness
claim; scheduled distillation, review and approval at scale, application
feedback and credibility decay went unexercised. One recorded lesson was
retrieved with its tier.

\subsection{E5: Content-layer configuration comparison}
\label{sec:ladder-opt}
The routing corpus was also used to compare eight evidence configurations.
A0 is the tiered baseline; A0f uses expanded document windows, A0h additionally
unions vector chunk windows, PD uses a navigation-free precise per-document
search, A1b a scoped flat union and reference-mode judge, A2 skips that judge,
A3 uses nominally name-only cards, and read-all scans every stored part.
A0h, PD and A1b differ in multiple components, so their differences cannot be
assigned to navigation alone. A3 may retrieve outside the evaluation scope.

In the 11 November aggregate, gold strings were contained for A0/A0f/A0h/PD/
A1b/A2/A3/read-all in 37/45/52/46/39/25/33/72 of 72 questions, respectively
(Figure~\ref{fig:ladder}); these are evidence-string proxies, not answer
correctness. Supplementary Tables S12--S13 give the per-arm intervals and
paired records. The A0h Wilson 95\% interval is [0.610, 0.812]. In exploratory,
unadjusted paired records, A0h contained the string while A1b did not on 13
questions and the reverse occurred on none (two-sided exact McNemar
$p=0.00024$). A0h led PD on seven questions while PD led on one
($p=0.0703$). The no-gate arm A2, which skips the judge on the tiered route,
contained the string for 25/72 against 37/72 for A0 (uncorrected exact McNemar
$p=0.0227$; the later record repeats 25/72 exactly), the arm contrast that
bears on the evidence gate. These records do not show a statistically resolved advantage
over PD and are not unbiased deployment estimates: A0, A1b and A3 can use
the gold domain to resolve judge references, with unknown frequency and impact.
As an external content-based reference point on the same libraries, a separate
40-question floor probe measured a gold-document hit rate of 60.0\% for a
lexical BM25 baseline; it uses different questions and metrics, so its number
is not a paired comparison.

A later same-question record yields 39/44/51/45/39/25 for the first six arms,
respectively. Per-question agreement with the frozen record is 94.4\% for A0,
98.6\% for A0f, A0h and PD, and 100\% for A1b and A2; the same question set
is reused and complete execution provenance is unavailable. In a separately
added later A3 record, containment is 25/72, with 13 missing top candidates
and 23 empty packs; their causes are not established. Later read-all is
incomplete. The frozen read-all is a mechanical substring scan, not an agent
that reads and answers. A0h recorded 309 KB and 13.8 s per question versus
24.3 MB and 253.7 s for that scan; the recorded-byte definitions differ and
may include overlap.

Among A0's 35 misses, A0f, A0h and PD contain the string in 8, 15 and 9 cases,
respectively. These sets overlap and are not additive. A post-hoc search
classified five first matching string locations outside a nominal window and
30 inside; the chosen matching file need not be the gold document or a
candidate actually read, so these categories do not identify the causes of
misses (Figure~\ref{fig:miss-overlap}).

\begin{figure*}[t]
\centering
\includegraphics[width=\textwidth]{fig_ladder_optimization}
\caption{Content-layer configurations on 72 in-domain questions.
(a)~Frozen gold-string containment with Wilson 95\% intervals (blue circles)
and later records on the same questions (orange triangles). The frozen
aggregate combines dates; A3 was added separately to the later record and
read-all was incomplete. The dashed line marks frozen A0. Read-all is a
mechanical string scan, not an answering agent. A1b may use the gold domain
for reference resolution; A0h and PD differ in both routing and chunk union.
(b)~Frozen-aggregate mean tool calls and trace-recorded kilobytes per
question; bytes can overlap. Paired comparisons are reported in the text.}
\label{fig:ladder}
\end{figure*}

\begin{figure}[t]
\centering
\includegraphics[width=\columnwidth]{fig_miss_overlap}
\caption{Overlap among configurations for the 35 questions where frozen A0
lacked the gold string. Each bar is a mutually exclusive outcome across A0f,
A0h and PD; the counts are not additive configuration gains. These are
substring results, not identified causes of failure.}
\label{fig:miss-overlap}
\end{figure}
